\documentclass[12pt,a4paper]{article}

% ── Pacotes ────────────────────────────────────────────────────────────────────
\usepackage[utf8]{inputenc}
\usepackage[T1]{fontenc}
\usepackage[english]{babel}
\usepackage{amsmath,amssymb,amsthm}
\usepackage{mathrsfs}
\usepackage{geometry}
\usepackage{xcolor}
\usepackage{hyperref}
\usepackage{booktabs}
\usepackage{array}
\usepackage{enumitem}

\geometry{
	a4paper,
	top=2.5cm, bottom=2.5cm,
	left=3cm, right=3cm
}

\hypersetup{
	colorlinks=true,
	linkcolor=blue,
	citecolor=blue,
	urlcolor=blue,
	pdftitle={Reformulacao Espectral do Motor de Heranca Estrutural},
	pdfauthor={Tiago Bandeira}
}

% ── Ambientes de Teorema ───────────────────────────────────────────────────────
\newtheorem{theorem}{Teorema}[section]
\newtheorem{lemma}[theorem]{Lema}
\newtheorem{proposition}[theorem]{Proposição}
\newtheorem{corollary}[theorem]{Corolário}

\theoremstyle{definition}
\newtheorem{definition}[theorem]{Definição}
\newtheorem{example}[theorem]{Exemplo}
\newtheorem{remark}[theorem]{Observação}

% ── Atalhos matemáticos ────────────────────────────────────────────────────────
\newcommand{\N}{\mathbb{N}}
\newcommand{\Z}{\mathbb{Z}}
\newcommand{\R}{\mathbb{R}}
\newcommand{\C}{\mathbb{C}}
\newcommand{\T}{\mathbb{T}}
\newcommand{\PP}{\mathcal{P}}          % conjunto dos primos
\newcommand{\F}{\mathcal{F}}           % fita-dobra
\newcommand{\G}{\mathcal{G}}           % grade
\newcommand{\LL}{\mathcal{L}}
\newcommand{\HRm}{\mathrm{HR}^{-}}
\newcommand{\HRp}{\mathrm{HR}^{+}}
\newcommand{\Lam}{\Lambda}             % von Mangoldt
\newcommand{\e}[1]{e^{2\pi i (#1)}}
\newcommand{\norm}[1]{\left\| #1 \right\|}
\newcommand{\abs}[1]{\left| #1 \right|}
\newcommand{\ip}[2]{\langle #1,\, #2 \rangle}
\DeclareMathOperator{\Spec}{Spec}

% ── Título ─────────────────────────────────────────────────────────────────────
\title{%
	\textbf{Reformulação Espectral do Motor de Herança Estrutural}\\[0.4em]
	\large Uma Abordagem por Sistemas Dinâmicos à Conjectura de Goldbach
}
\author{Tiago Bandeira\\[0.3em]
	\small\texttt{tiagobandeira.goldbach2025@gmail.com}
}
\date{Maio de 2026 \\ \vspace{0.3em}
	\small Preprint — Artigo 6 da série sobre a Conjectura de Goldbach}

% ═══════════════════════════════════════════════════════════════════════════════
\begin{document}
	
	\maketitle
	\thispagestyle{empty}
	
	% ── Nota do Autor ─────────────────────────────────────────────────────────────
	\begin{center}
		\fbox{\parbox{0.88\textwidth}{
				\textbf{Nota do Autor.} Este artigo não reivindica uma demonstração da
				Conjectura de Goldbach. Seu objetivo é reformular a questão central —
				a Hipótese Restritiva Fraca $\HRm$ — na linguagem de sistemas dinâmicos
				e análise espectral, e provar incondicionalmente que toda a interferência
				espectral das frequências não nulas se cancela (via teorema de
				Green-Tao-Ziegler \cite{greentaoziegler}), reduzindo o problema a uma
				única componente de frequência zero. Os resultados incondicionais são
				explicitamente identificados; os condicionais são declarados como tais.
				Comentários e críticas matemáticas são bem-vindos.
		}}
	\end{center}
	\vspace{0.5em}
	
	% ── Abstract ──────────────────────────────────────────────────────────────────
	\begin{abstract}
		Reformulamos o problema central do Motor de Herança Estrutural
		\cite{motor} na linguagem de sistemas dinâmicos e análise espectral.
		A Hipótese Restritiva Fraca ($\HRm$) — existência de um eixo com par
		de primos em alguma janela $W_i$ da varredura do scanner — é
		traduzida, de forma incondicional, para a condição
		$\widehat{F}(0) > 0$: positividade da componente de frequência zero
		da função indicadora de colisão
		$F(t) = \mathbf{1}_{\PP}(E_1(t))\cdot\mathbf{1}_{\PP}(E_2(t))$
		no sistema dinâmico de medida finito $(X_N, \mu_N, T)$, onde $t$
		indexa as diagonais da fita (pares que somam o par alvo $2M^+$).
		
		Demonstramos incondicionalmente que o operador de Koopman $U_T$
		possui espectro discreto puro (raízes $K$-ésimas da unidade,
		$K = N-1$), que $T$ é ergódico com entropia de Kolmogorov-Sinai
		nula, e que $\HRm \Leftrightarrow \widehat{F}(0) > 0$.
		Adicionalmente, provamos de forma incondicional — via o teorema de
		Green-Tao-Ziegler \cite{greentaoziegler} sobre ortogonalidade da
		função de von Mangoldt a nilsequências — que
		$\widehat{w}(j) \to 0$ para todo $j \neq 0$.
		O problema se reduz assim, incondicionalmente, a uma única
		componente espectral: a positividade de $\widehat{w}(0)$, que é
		equivalente a $\HRm$ e permanece em aberto.
		
		Discutimos as razões precisas pelas quais os métodos clássicos —
		método do círculo de Hardy-Littlewood-Vinogradov e crivo linear de
		Rosser-Iwaniec — colapsam nesta configuração, e delineamos como a
		formulação espectral reposiciona a barreira analítica remanescente em
		termos acessíveis às técnicas modernas.
	\end{abstract}
	
	\tableofcontents
	\newpage
	
	% ═══════════════════════════════════════════════════════════════════════════════
	\section{Introdução}
	\label{sec:intro}
	
	\subsection{A série e o contexto}
	
	Este artigo integra a série \cite{goldbach1,grade,ancora,crivo_b,motor}
	dedicada a construir um mecanismo dinâmico de cobertura dos números pares
	pela Conjectura de Goldbach. O Motor de Herança Estrutural \cite{motor}
	reduz o problema global
	\begin{quote}
		\emph{``Todo par $2M > 2$ é soma de dois primos''}
	\end{quote}
	à questão local, repetida para cada iteração do motor:
	\begin{quote}
		\emph{``Existe algum eixo da janela $W_i$, para algum $i$ na
			varredura do scanner, tal que ambos os extremos $E_1$ e $E_2$
			desse eixo são primos?''}
	\end{quote}
	Esta é precisamente $\HRm$. A camada algébrica do motor —
	construção da Fita-Dobra, acoplamento com a Grade, preenchimento
	de $W_i$ pelo scanner — foi demonstrada incondicionalmente em
	\cite{motor}. A única lacuna remanescente é a primalidade
	simultânea dos extremos de algum eixo de $W_i$.
	
	O presente artigo não resolve essa lacuna. Seu propósito é
	reescrevê-la na linguagem de sistemas dinâmicos e análise
	harmônica, tornando explícita a natureza da obstrução e alinhando
	o problema com a escola moderna de teoria analítica dos números.
	
	\subsection{Por que os métodos clássicos falham}
	
	As três principais tentativas de provar $\HRm$ diretamente pela
	teoria analítica dos números clássica colapsam por razões
	estruturais precisas, detalhadas na Seção~\ref{sec:crivo}:
	
	\begin{enumerate}[label=(\roman*)]
		\item \textbf{Método do círculo (Hardy-Littlewood-Vinogradov).} A
		separação de escalas — comprimento de progressão
		$L \sim (\log M^-)^2$ versus alvo $2M^- \gg L$ — destrói o
		termo principal da integral de Fourier quando $M \to \infty$.
		
		\item \textbf{Crivo linear de Rosser-Iwaniec.} O erro acumulado
		$\mathcal{R}$ cresce como $(\log M^-)^{3-\varepsilon}$, dominando
		o termo principal de ordem $(\log M^-)^2 / (\log\log M^-)^2$.
		
		\item \textbf{Indução algébrica pura.} A primalidade é
		multiplicativa; a fita-dobra opera sob leis aditivas. Transportar
		primalidade passo a passo cai em circularidade indutiva.
	\end{enumerate}
	
	O diagnóstico comum é claro: \emph{ferramentas estáticas de ponto
		fixo não resolvem um sistema dinâmico em evolução}. A reformulação
	proposta analisa o fluxo orbital global em vez de instâncias
	isoladas.
	
	\subsection{Mudança de perspectiva: da contagem ao fluxo}
	
	A filosofia central é a seguinte transição de linguagem:
	
	\begin{center}
		\renewcommand{\arraystretch}{1.4}
		\begin{tabular}{lll}
			\toprule
			\textbf{Dimensão} & \textbf{Abordagem Clássica} & \textbf{Abordagem Espectral} \\
			\midrule
			Análise & Pontual (um $2M$ de cada vez) & Global (fluxo sobre $X_N$) \\
			Matemática & Aritmética de contagem & Geometria de medida \\
			O motor & Organizador geométrico & Operador espectral (Koopman) \\
			A barreira & Paridade e acúmulo de $\mathcal{R}$ & Freq.\ zero de $\widehat{w}$ (GTZ isola) \\
			\bottomrule
		\end{tabular}
	\end{center}
	
	Esta transição sintoniza o Motor de Herança com a escola moderna de
	Teoria Analítica dos Números (Green-Tao-Gowers \cite{greentao},
	Green-Tao-Ziegler \cite{greentaoziegler}, Sarnak \cite{sarnak}),
	na qual a pseudo-aleatoriedade dos primos é explorada como
	ferramenta, não como obstáculo.
	
	\begin{remark}[O que a reformulação oferece e o que não oferece]
		\label{rem:honestidade}
		A equivalência $\HRm \Leftrightarrow \widehat{F}(0) > 0$
		(Teorema~\ref{thm:reformulacao}) é uma tautologia precisa, não
		uma prova. Ela não reduz a dificuldade de Goldbach; ela a
		\emph{localiza} numa obstrução espectral única e bem definida.
		O valor está em que essa obstrução é exatamente o tipo de objeto
		para o qual existem ferramentas modernas — nilsequências, normas
		de Gowers, princípio de Sarnak — que ainda não estavam disponíveis
		nas formulações clássicas.
	\end{remark}
	
	\subsection{Organização do artigo}
	
	A Seção~\ref{sec:fundamentos} recapitula os fundamentos algébricos
	do motor: Fita-Dobra, Grade com percurso zigzag, Janela $W_i$ com
	seus quatro eixos de simetria, e o Scanner com mecanismo de pivôs.
	A Seção~\ref{sec:dinamica} constrói o sistema dinâmico
	$(X_N, \mu_N, T)$ e prova o Teorema Espectral (incondicional).
	A Seção~\ref{sec:colisao} define a função de colisão $F$ e
	estabelece $\HRm \Leftrightarrow \widehat{F}(0) > 0$
	(incondicional). A Seção~\ref{sec:obstrucao} descreve a obstrução
	espectral e sua conexão com Sarnak e Gowers. A
	Seção~\ref{sec:crivo} detalha o colapso dos métodos clássicos. A
	Seção~\ref{sec:conclusao} sintetiza o alcance e aponta direções
	futuras.
	
	% ═══════════════════════════════════════════════════════════════════════════════
	\section{Fundamentos Algébricos do Motor}
	\label{sec:fundamentos}
	
	Nesta seção apresentamos os objetos e resultados incondicionais de
	\cite{motor} necessários para a formulação espectral. A distinção
	central que organiza toda a estrutura é:
	
	\begin{center}
		\renewcommand{\arraystretch}{1.4}
		\begin{tabular}{lll}
			\toprule
			\textbf{Conceito} & \textbf{Definição} & \textbf{Quem certifica} \\
			\midrule
			Par base $2M^-$ & Soma das \emph{colunas} da fita: $a_k + b_k = 2M^-$
			& Goldbach conhecido \\
			Par alvo $2M^+$ & Soma das \emph{diagonais} da fita:
			$a_{k+1} + b_k = 2M^+$
			& $\HRm$ certifica via $W_i$ \\
			\bottomrule
		\end{tabular}
	\end{center}
	
	Todos os resultados desta seção são independentes de primalidade.
	
	\subsection{A Fita-Dobra}
	
	\begin{definition}[Fita-Dobra $\F(2 \times N)$]
		\label{def:fita}
		Fixado o par base $2M^-$ com $M^- \geq 2$, a Fita-Dobra
		$\F(2 \times N)$ é a matriz $2 \times N$ com $N = M^-$ entradas por
		linha, definida por
		\[
		a_k = 2k - 1,\qquad b_k = 2M^- - (2k-1),\qquad k = 1,\ldots,N.
		\]
		A Linha~1 percorre os ímpares crescentes $1, 3, 5, \ldots, 2N-1$
		da esquerda para a direita. A Linha~2 percorre os complementares
		$2M^{-}-1, 2M^{-}-3, \ldots, 1$ da direita para a esquerda.
	\end{definition}
	
	\begin{proposition}[Soma constante nas colunas — par base]
		\label{prop:coluna}
		Para toda coluna $k \in \{1, \ldots, N\}$:
		\[
		a_k + b_k = 2M^-.
		\]
		Esta soma é constante, independente de $k$ e independente de
		primalidade. O par base $2M^-$ é certificado ao se encontrar uma
		coluna $k$ com $a_k$ e $b_k$ ambos primos — o motor parte desse
		fato como dado.
	\end{proposition}
	
	\begin{proposition}[Soma constante nas diagonais — par alvo]
		\label{prop:diagonal}
		Para todo $k \in \{1, \ldots, N-1\}$, o par diagonal
		$(a_{k+1},\, b_k)$ satisfaz
		\[
		a_{k+1} + b_k
		= (2k+1) + \bigl(2M^- - (2k-1)\bigr)
		= 2M^- + 2
		=: 2M^+.
		\]
		Cada par diagonal soma exatamente $2M^+ = 2M^- + 2$,
		independente de primalidade. Esses $N-1$ pares diagonais são o
		objeto capturado pelo Scanner e armazenado nos eixos de $W_i$.
	\end{proposition}
	
	\begin{remark}[A fita representa dois pares consecutivos]
		A Fita-Dobra representa simultaneamente o par base $2M^-$
		(soma das colunas) e o par alvo $2M^+$ (soma das diagonais).
		O par alvo de uma iteração torna-se o par base da iteração
		seguinte, formando a cadeia de herança estrutural.
	\end{remark}
	
	\subsection{A Grade com Percurso Zigzag}
	
	\begin{definition}[Grade $\G(3 \times C)$ com percurso zigzag]
		\label{def:grade}
		A Grade $\G(3 \times C)$ é uma matriz de $3$ linhas e $C$ colunas
		que contém os mesmos $3C$ ímpares da Fita-Dobra, organizados em
		percurso zigzag de três linhas:
		\begin{itemize}
			\item Linha~1: percorre os ímpares crescentes da
			\emph{esquerda para a direita} ($\rightarrow$).
			\item Linha~2: percorre os ímpares seguintes da
			\emph{direita para a esquerda} ($\leftarrow$).
			\item Linha~3: percorre os ímpares seguintes da
			\emph{esquerda para a direita} ($\rightarrow$).
		\end{itemize}
		O percurso zigzag garante que as diagonais geométricas da grade
		correspondam às diagonais da fita.
	\end{definition}
	
	\begin{definition}[Condição de acoplamento]
		\label{def:acoplamento}
		Para que Grade e Fita representem exatamente o mesmo conjunto de
		ímpares, a condição de acoplamento deve ser satisfeita:
		\[
		3C = 2N.
		\]
		Quando o total de ímpares não é divisível por~$3$, duplicações
		centrais são introduzidas para satisfazer a condição sem alterar
		a cardinalidade efetiva do sistema.
	\end{definition}
	
	\begin{example}[Grade $\G(3 \times 8)$, par base $2M^- = 48$,
		par alvo $2M^+ = 50$]
		Acoplamento: $3 \times 8 = 2 \times 12 = 24$ células. $\checkmark$
		\[
		\begin{array}{cccccccc}
			1  & 3  & 5  & 7  & 9  & 11 & 13 & 15 \\
			31 & 29 & 27 & 25 & 23 & 21 & 19 & 17 \\
			33 & 35 & 37 & 39 & 41 & 43 & 45 & 47
		\end{array}
		\]
	\end{example}
	
	\subsection{A Janela \texorpdfstring{$W_i$}{Wi} — Quatro Eixos de Simetria}
	
	\begin{definition}[Janela $W_i$]
		\label{def:wi}
		A Janela $W_i$ é uma matriz $3 \times 3$ fixa que armazena, nos
		seus quatro eixos de simetria, até quatro pares diagonais da
		fita, cada um somando o par alvo $2M^+$. $W_i$ não é deslizante:
		ela é preenchida pelo Scanner e permanece como referência
		estrutural.
		
		Uma matriz $3 \times 3$ possui exatamente quatro linhas de simetria
		que passam pelo centro, cada uma conectando dois extremos opostos:
		\begin{center}
			\renewcommand{\arraystretch}{1.3}
			\begin{tabular}{lll}
				\toprule
				\textbf{Eixo} & \textbf{Entrada $E_1$} & \textbf{Entrada $E_2$ (oposto)} \\
				\midrule
				$D_1$ (diagonal principal) & canto $[0,0]$ & canto $[2,2]$ \\
				$D_2$ (diagonal secundária) & canto $[0,2]$ & canto $[2,0]$ \\
				$C_v$ (coluna central) & topo $[0,1]$ & base $[2,1]$ \\
				$L_c$ (linha central) & esquerda $[1,0]$ & direita $[1,2]$ \\
				\bottomrule
			\end{tabular}
		\end{center}
		Em cada eixo: o elemento menor do par é posicionado na entrada
		$E_1$, o maior em $E_2$. A posição central $[1,1]$ armazena o
		\textbf{pivô} — o elemento central da varredura do Scanner.
		A estrutura visual de $W_i$ é:
		\[
		\begin{pmatrix}
			E_1(D_1) & E_1(C_v) & E_1(D_2) \\
			E_1(L_c) & \text{PIV\^O} & E_2(L_c) \\
			E_2(D_2) & E_2(C_v) & E_2(D_1)
		\end{pmatrix}
		\]
	\end{definition}
	
	\begin{proposition}[Soma constante nos eixos de $W_i$]
		\label{prop:eixos}
		Por construção da Fita-Dobra (Proposição~\ref{prop:diagonal}),
		cada par $(E_1, E_2)$ capturado pelo Scanner e armazenado num
		eixo de $W_i$ satisfaz automaticamente
		\[
		E_1 + E_2 = 2M^+,
		\]
		independente de os elementos serem ou não primos. Esta é a
		garantia algébrica incondicional do motor.
	\end{proposition}
	
	\begin{example}[$W_i$ com par base $2M^- = 48$, par alvo $2M^+ = 50$]
		\[
		\begin{pmatrix}
			7  &  9 & 11 \\
			37 & 25 & 13 \\
			39 & 41 & 43
		\end{pmatrix}
		\]
		$D_1$: $7 + 43 = 50$ $\checkmark$\quad
		$C_v$: $9 + 41 = 50$ $\checkmark$\quad
		$D_2$: $11 + 39 = 50$ $\checkmark$\quad
		$L_c$: $13 + 37 = 50$ $\checkmark$
		
		$\HRm$ ativada em $D_1$: $7$ primo $\checkmark$,
		$43$ primo $\checkmark$ $\Rightarrow$ $50 = 7 + 43$.
	\end{example}
	
	\subsection{O Scanner — Mecanismo dos Pivôs}
	
	\begin{definition}[Scanner com pivôs]
		\label{def:scanner}
		O Scanner percorre a Grade e preenche os eixos de $W_i$ com as
		diagonais da Fita. Dois pivôs marcham das extremidades da lista
		de ímpares em direção ao centro:
		\begin{enumerate}
			\item Posicionar o pivô esquerdo na \emph{segunda} posição da
			lista (o elemento~$1$ é descartado, pois não forma
			diagonal secundária válida: não existe $a$ na grade com
			$1 + a = 2M^+$).
			\item Posicionar o pivô direito na extremidade final da lista.
			\item Registrar o par $(E_{\mathrm{esq}}, E_{\mathrm{dir}})$
			num eixo de $W_i$. O elemento central da lista é o
			pivô do sensor (posição $[1,1]$).
			\item Avançar o pivô esquerdo ($+1$) e recuar o pivô direito
			($-1$).
			\item Repetir até os pivôs se encontrarem. Cada iteração
			preenche um eixo de $W_i$.
			\item Quando $W_i$ tiver os quatro eixos preenchidos,
			registrá-la e iniciar nova $W_i$ para o bloco seguinte.
		\end{enumerate}
		O número de janelas necessárias para cobrir todas as $N-1$
		diagonais da fita é $\lceil (N-1)/4 \rceil$.
	\end{definition}
	
	\subsection{As Hipóteses Restritivas}
	
	\begin{definition}[Hipóteses Restritivas]
		\label{def:HR}
		Fixados $N$ e a varredura do Scanner sobre a Grade, sejam
		$(E_1, E_2)$ os extremos de um eixo de $W_i$, com
		$E_1 + E_2 = 2M^+$ por construção geométrica
		(Proposição~\ref{prop:eixos}).
		
		\begin{itemize}
			\item \textbf{$\HRm$ (Fraca):} Existe pelo menos um eixo de
			$W_i$, para alguma janela $i$ na varredura, tal que
			$E_1 \in \PP$ e $E_2 \in \PP$. O par $(E_1, E_2)$ é então
			um certificado de Goldbach para o par alvo $2M^+$.
			\item \textbf{$\HRp$ (Forte):} Existe pelo menos um eixo de
			$W_i$ tal que a tripla completa
			$(E_1, \mathrm{Piv\hat{o}}, E_2)$ é inteiramente formada
			por primos. $\HRp$ implica $\HRm$.
		\end{itemize}
	\end{definition}
	
	Sob $\HRm$, o Teorema do Motor \cite[Thm.~9.1]{motor} garante que
	a cadeia
	\[
	2M_0^- \;\to\; 2M_0^+ \;\to\; 2M_0^+ + 2 \;\to\; \cdots
	\]
	cobre todos os pares $\geq 2M_0^-$ como somas de dois primos.
	A lacuna central — a única hipótese não provada — é precisamente
	$\HRm$.
	
	\begin{remark}[Separação entre camadas algébrica e aritmética]
		A arquitetura do motor separa claramente duas camadas. A
		\emph{camada algébrica} — construção da fita, acoplamento com
		a grade, estrutura de $W_i$, preenchimento pelo Scanner — é
		inteiramente incondicional: a existência geométrica dos pares e
		sua cobertura pelos eixos de $W_i$ independem de qualquer
		hipótese sobre primos. A \emph{camada aritmética} — a
		primalidade simultânea dos extremos de algum eixo — é o único
		ingrediente não provado, capturado por $\HRm$.
	\end{remark}
	
	% ═══════════════════════════════════════════════════════════════════════════════
	\section{O Sistema Dinâmico de Medida Finito}
	\label{sec:dinamica}
	
	\subsection{Construção do sistema}
	
	A varredura do Scanner sobre a Grade produz exatamente $N-1$ pares
	diagonais $(E_1(t), E_2(t))$, indexados por $t \in \{1,\ldots,N-1\}$,
	cada um satisfazendo $E_1(t)+E_2(t)=2M^+$ incondicionalmente. O
	sistema dinâmico é construído sobre esse conjunto de estados.
	
	\begin{definition}[Espaço de estados $X_N$]
		\label{def:estados}
		O espaço de estados do motor é o conjunto finito de índices
		diagonais
		\[
		X_N := \{t \in \Z : 1 \leq t \leq N-1\}.
		\]
		Cada estado $t \in X_N$ corresponde biunivocamente ao par
		diagonal $(E_1(t), E_2(t))$ capturado pelo Scanner, com soma
		$E_1(t)+E_2(t)=2M^+$ garantida geometricamente pela
		Proposição~\ref{prop:diagonal}.
	\end{definition}
	
	\begin{definition}[Medida de probabilidade uniforme $\mu_N$]
		Definimos a medida de probabilidade sobre a $\sigma$-álgebra
		discreta de $X_N$ por
		\[
		\mu_N(\{t\}) = \frac{1}{N-1},\qquad \forall\, t \in X_N.
		\]
	\end{definition}
	
	\begin{definition}[Transformação dinâmica $T$]
		A transformação $T: X_N \to X_N$ é a translação cíclica unitária
		\[
		T(t) = t + 1 \pmod{N-1}.
		\]
		$T$ percorre ciclicamente todos os índices diagonais,
		equivalendo a observar o par seguinte na enumeração do Scanner.
	\end{definition}
	
	\begin{lemma}[Preservação de medida e ergodicidade]
		\label{lem:ergodico}
		A transformação $T$ é um automorfismo de medida sobre
		$(X_N, \mu_N)$. Além disso, $(X_N, \mu_N, T)$ é ergódico e
		possui entropia métrica de Kolmogorov-Sinai nula:
		\[
		\mu_N(T^{-1}(B)) = \mu_N(B)\;\;\forall B \subset X_N,
		\qquad h(T) = 0.
		\]
	\end{lemma}
	
	\begin{proof}
		Como $T$ é uma bijeção de $X_N$ sobre si mesmo e $\mu_N$ é a
		medida de contagem normalizada, a preservação de medida é
		imediata. A ergodicidade decorre do fato de que $T$ forma um
		único ciclo de comprimento máximo $K = N-1$ sobre $X_N$: o
		único subconjunto $B$ com $T^{-1}(B) = B$ é $B = \emptyset$ ou
		$B = X_N$. A entropia nula é consequência do caráter periódico
		e determinístico de $T$.
	\end{proof}
	
	\subsection{O operador de Koopman e sua estrutura espectral}
	
	\begin{definition}[Espaço de Hilbert e produto interno]
		Seja $L^2(X_N, \mu_N)$ o espaço de funções de valor complexo
		sobre $X_N$, com produto interno
		\[
		\ip{\psi}{\phi} = \frac{1}{N-1}\sum_{t=1}^{N-1}
		\psi(t)\,\overline{\phi(t)}.
		\]
	\end{definition}
	
	\begin{definition}[Operador de Koopman $U_T$]
		\label{def:koopman}
		O operador de Koopman $U_T: L^2(X_N) \to L^2(X_N)$ é definido
		por
		\[
		(U_T\psi)(t) = \psi(T(t)).
		\]
		Como $T$ preserva $\mu_N$, $U_T$ é unitário: $U_T^* = U_T^{-1}$.
		Geometricamente, $U_T$ desloca o índice de observação ao longo
		da sequência de pares diagonais do Scanner.
	\end{definition}
	
	\begin{proposition}[Representação matricial circulante]
		\label{prop:circulante}
		Na base canônica de $L^2(X_N)$, o operador $U_T$ é representado
		pela matriz de permutação circulante
		$M \in \R^{K\times K}$ (com $K = N-1$):
		\[
		M = \begin{pmatrix}
			0 & 1 & 0 & \cdots & 0 \\
			0 & 0 & 1 & \cdots & 0 \\
			\vdots & \vdots & \vdots & \ddots & \vdots \\
			0 & 0 & 0 & \cdots & 1 \\
			1 & 0 & 0 & \cdots & 0
		\end{pmatrix}.
		\]
	\end{proposition}
	
	\begin{theorem}[Espectro discreto do motor — incondicional]
		\label{thm:espectro}
		O espectro de $U_T$ é discreto puro. Os autovalores e as
		autofunções associadas são:
		\[
		\lambda_j = \exp\!\left(\frac{2\pi i\, j}{K}\right),\qquad
		\phi_j(t) = \exp\!\left(\frac{2\pi i\, j\, t}{K}\right),\qquad
		j,t \in \{0,\ldots,K-1\}.
		\]
		Em particular, $\phi_0 \equiv 1$ é a autofunção constante
		associada ao autovalor trivial $\lambda_0 = 1$.
	\end{theorem}
	
	\begin{proof}
		Como $M$ é uma matriz circulante de permutação, satisfaz
		$M^K = I$, donde seu polinômio mínimo divide $x^K - 1$. As
		raízes de $x^K - 1$ são exatamente $\lambda_j = e^{2\pi ij/K}$,
		$j = 0,\ldots,K-1$. As autofunções $\phi_j$ são os harmônicos
		de Fourier discretos sobre $\Z_K$; sua ortogonalidade decorre
		da relação de ortogonalidade de caracteres:
		\[
		\ip{\phi_j}{\phi_\ell} = \frac{1}{K}\sum_{t=0}^{K-1}
		e^{2\pi i(j-\ell)t/K} = \delta_{j\ell}.\qedhere
		\]
	\end{proof}
	
	\begin{remark}[Significado do Teorema~\ref{thm:espectro}]
		O resultado demonstra, sem qualquer hipótese sobre primos, que
		a dinâmica do motor é geometricamente equivalente a um
		\emph{oscilador harmônico discreto} de entropia nula. Esta
		propriedade é o ponto de entrada para a Conjectura de Sarnak:
		sistemas de entropia zero têm comportamento ``simples'' o
		suficiente para que a pseudo-aleatoriedade dos primos se
		manifeste como descorrelação espectral.
	\end{remark}
	
	% ═══════════════════════════════════════════════════════════════════════════════
	\section{Função de Colisão e Reformulação Espectral de
		\texorpdfstring{$\HRm$}{HR-}}
	\label{sec:colisao}
	
	\subsection{A função indicadora de colisão}
	
	\begin{definition}[Função de colisão $F$]
		\label{def:colisao}
		A observável aritmética $F \in L^2(X_N)$ que registra colisões
		primas nos pares diagonais do motor é definida por
		\[
		F(t) := \mathbf{1}_{\PP}(E_1(t))\cdot\mathbf{1}_{\PP}(E_2(t)),
		\]
		onde $\mathbf{1}_{\PP}(x) = 1$ se $x$ é primo e $0$ caso
		contrário, e $(E_1(t), E_2(t))$ é o $t$-ésimo par diagonal
		capturado pelo Scanner com $E_1(t)+E_2(t)=2M^+$.
	\end{definition}
	
	\begin{definition}[Função de peso de von Mangoldt]
		\label{def:mangoldt}
		A versão suavizada de $F$ via a função de von Mangoldt é
		\[
		w(t) := \Lambda(E_1(t))\cdot\Lambda(E_2(t)),
		\]
		onde $\Lambda(n) = \log p$ se $n = p^k$ e $0$ caso contrário.
		A passagem de $w$ para $F$ introduz erro $O(\sqrt{K}\log K)$
		proveniente das potências de primos com $k \geq 2$.
	\end{definition}
	
	\subsection{Equações de indexação dos pares diagonais}
	
	\begin{remark}[Variação linear dos extremos diagonais]
		\label{rem:movimento}
		Por construção da Fita-Dobra (Definição~\ref{def:fita}), os
		extremos do $t$-ésimo par diagonal variam linearmente:
		\[
		E_1(t) = E_1(1) + 2(t-1),\qquad
		E_2(t) = E_2(1) - 2(t-1),
		\]
		de modo que $E_1(t) + E_2(t) = E_1(1) + E_2(1) = 2M^+$ é
		constante para todo $t \in \{1,\ldots,N-1\}$.
	\end{remark}
	
	\subsection{Decomposição espectral e reformulação de
		\texorpdfstring{$\HRm$}{HR-}}
	
	\begin{definition}[Transformada de Fourier de $F$]
		A projeção espectral de $F$ sobre os autovalores do motor é
		\[
		\widehat{F}(j) := \ip{F}{\phi_j} = \frac{1}{K}\sum_{t=1}^{K}
		F(t)\,e^{-2\pi i j t/K}.
		\]
	\end{definition}
	
	\begin{theorem}[Reformulação espectral de $\HRm$ — incondicional]
		\label{thm:reformulacao}
		A Hipótese Restritiva Fraca $\HRm$ para o par alvo $2M^+$ é
		equivalente a qualquer das seguintes condições:
		\begin{enumerate}[label=(\alph*)]
			\item $F(t_0) \neq 0$ para algum $t_0 \in X_N$.
			\item $\displaystyle\sum_{t=1}^{K} F(t) \geq 1$.
			\item $\widehat{F}(0) = \ip{F}{\mathbf{1}} > 0$.
		\end{enumerate}
	\end{theorem}
	
	\begin{proof}
		A equivalência (a)$\Leftrightarrow$(b) é imediata pois
		$F(t) \in \{0,1\}$. Para (b)$\Leftrightarrow$(c): como
		$\phi_0 \equiv 1$ e $\mu_N$ é uniforme,
		\[
		\widehat{F}(0) = \frac{1}{K}\sum_{t=1}^K F(t) \geq 0,
		\]
		com igualdade se e somente se $F \equiv 0$. Portanto
		$\widehat{F}(0) > 0 \Leftrightarrow \sum F(t) \geq 1$.
	\end{proof}
	
	\begin{remark}[Natureza da reformulação]
		\label{rem:tautologia}
		O Teorema~\ref{thm:reformulacao} é uma equivalência, não uma
		prova. A condição $\widehat{F}(0) > 0$ é precisamente uma
		reescrita de $\HRm$ na linguagem de análise harmônica. O ganho
		não é lógico, mas \emph{linguístico}: a obstrução é agora um
		objeto espectral para o qual existem técnicas modernas —
		descorrelação de Sarnak, normas de Gowers, nilsequências — que
		não têm análogo direto na formulação original.
	\end{remark}
	
	% ═══════════════════════════════════════════════════════════════════════════════
	\section{A Obstrução Espectral e Conexões com Ferramentas Modernas}
	\label{sec:obstrucao}
	
	\subsection{Pseudo-aleatoriedade e uniformidade de Gowers}
	
	\begin{definition}[Uniformidade de Gowers $U^2$]
		Uma sequência $a: \{1,\ldots,K\} \to \C$ é
		\emph{$U^2$-uniforme} se
		\[
		\norm{a}_{U^2(K)}^4 := \frac{1}{K^3}\sum_{n,h_1,h_2}
		a(n)\overline{a(n+h_1)}\overline{a(n+h_2)}a(n+h_1+h_2) = o(1)
		\quad (K \to \infty).
		\]
	\end{definition}
	
	Um resultado central na combinatória aditiva moderna é que a
	função de von Mangoldt normalizada é $U^2$-uniforme ao longo de
	progressões aritméticas (Green-Tao \cite{greentao}, Gowers
	\cite{gowers}). Intuitivamente, isso significa que os primos não
	exibem correlações a dois pontos em qualquer progressão aritmética
	fixa — propriedade que, na linguagem do motor, se traduz como
	descorrelação entre $w(t)$ e as frequências $j \neq 0$ do sistema.
	
	\subsection{Descorrelação Espectral via Nilsequências}
	
	A versão anterior deste resultado invocava a Conjectura de Sarnak
	\cite{sarnak} — sobre a ortogonalidade de $\mu(n)$ a sistemas de
	entropia zero — para obter descorrelação de $\Lambda$. Essa
	passagem de $\mu$ para $\Lambda$ é não trivial e torna o resultado
	condicional. O teorema de Green-Tao-Ziegler \cite{greentaoziegler}
	oferece uma rota incondicional mais direta.
	
	\begin{proposition}[Descorrelação espectral — incondicional]
		\label{prop:sarnak}
		Seja $(X_N, \mu_N, T)$ o sistema dinâmico do motor. Para todo
		$j \neq 0$, a autofunção
		\[
		\phi_j(t) = \exp\!\left(\frac{2\pi i\,j\,t}{K}\right)
		\]
		é uma nilsequência de grau $1$ (caractere linear sobre $\Z_K$).
		Pelo teorema de Green-Tao-Ziegler \cite{greentaoziegler} sobre
		sistemas lineares em primos — especificamente, sobre a
		ortogonalidade do produto $\Lambda(L_1(t))\Lambda(L_2(t))$ a
		nilsequências para pares de formas lineares não degenerados —
		a função de peso $w(t) = \Lambda(E_1(t))\,\Lambda(E_2(t))$
		satisfaz
		\[
		\lim_{K\to\infty}\widehat{w}(j) = 0,\qquad \forall\, j \neq 0.
		\]
	\end{proposition}
	
	\begin{proof}
		Consideramos o par de formas lineares
		\[
		L_1(t) = E_1(1) + 2(t-1),\qquad
		L_2(t) = E_2(1) - 2(t-1),
		\]
		com coeficientes $+2$ e $-2$ (Observação~\ref{rem:movimento}).
		Este par é \emph{não degenerado} — os coeficientes são
		linearmente independentes e a progressão é localmente solúvel
		— e satisfaz $L_1(t) + L_2(t) = 2M^+$ constante.
		
		Pelo teorema principal de Green-Tao-Ziegler
		\cite{greentaoziegler}, para qualquer nilsequência $\psi$ de
		grau finito,
		\[
		\frac{1}{K}\sum_{t=1}^{K} \Lambda(L_1(t))\,\Lambda(L_2(t))\,
		\psi(t) \to 0 \quad (K \to \infty).
		\]
		Para cada $j \neq 0$, a autofunção $\phi_j(t) = e^{2\pi i j t/K}$
		é uma nilsequência de grau $1$. Tomando
		$\psi(t) = \overline{\phi_j(t)}$, obtemos
		\[
		\widehat{w}(j) = \frac{1}{K}\sum_{t=1}^{K}
		\Lambda(E_1(t))\,\Lambda(E_2(t))\,\overline{\phi_j(t)}
		\to 0. \qedhere
		\]
	\end{proof}
	
	\begin{remark}[Relação com a Conjectura de Sarnak]
		\label{rem:sarnak}
		A Conjectura de Sarnak \cite{sarnak} — sobre $\mu(n)$ ortogonal
		a sistemas de entropia zero — é mais geral mas ainda não provada.
		O resultado acima não a utiliza: ele decorre diretamente do
		teorema de nilsequências de Green-Tao-Ziegler, que é
		incondicional. A Conjectura de Sarnak permanece como motivação
		conceitual para a abordagem, mas não como hipótese necessária.
	\end{remark}
	
	\begin{remark}[Nota sobre possível novidade na literatura]
		\label{rem:novidade}
		A aplicação do teorema de Green-Tao-Ziegler diretamente ao
		sistema de Koopman do Motor de Herança — concluindo descorrelação
		incondicional das componentes $j \neq 0$ — pode constituir uma
		observação nova. Não identificamos na literatura uma aplicação
		análoga de nilsequências a sistemas de cobertura de pares de
		Goldbach com soma fixa desta forma. Entretanto, a literatura em
		teoria ergódica e teoria analítica dos números é vasta, e uma
		verificação cuidadosa por especialistas é necessária antes de
		qualquer reivindicação de novidade.
	\end{remark}
	
	\subsection{A obstrução espectral e o que ela implica}
	
	A Proposição~\ref{prop:sarnak} estabelece que as componentes
	$\widehat{w}(j)$ para $j \neq 0$ tendem a zero. A componente
	$\widehat{w}(0)$, por sua vez, é a média orbital de $w$, e é
	precisamente o objeto cuja positividade equivale a $\HRm$.
	
	\begin{remark}[A barreira remanescente — formulação precisa]
		\label{rem:barreira}
		A lacuna analítica central pode ser formulada com precisão como:
		\[
		\text{provar incondicionalmente que}\quad
		\frac{1}{K}\sum_{t=1}^K \Lambda(E_1(t))\,\Lambda(E_2(t)) \gg
		\frac{1}{(\log K)^2},
		\]
		onde $E_1(t)+E_2(t) = 2M^+$ para todo $t$, por construção
		geométrica da Fita-Dobra. Esta estimativa é equivalente à
		própria Conjectura de Goldbach na linguagem da análise harmônica.
		A reformulação não simplifica essa dificuldade; ela a expressa
		num objeto — a componente de frequência zero de $w$ no sistema
		de Koopman — para o qual as técnicas de nilsequências e médias
		de Weyl têm tração natural.
	\end{remark}
	
	\begin{remark}[Conexão com a heurística de Hardy-Littlewood]
		\label{rem:HL}
		A Conjectura de Hardy-Littlewood \cite{hardy-littlewood} prediz
		que
		\[
		\widehat{w}(0) \sim \mathfrak{S}(2M^+)\cdot
		\frac{(\log M^+)^2}{(\log K)^2} > 0,
		\]
		onde $\mathfrak{S}(2M^+)$ é a série singular associada ao par
		alvo, sempre positiva.
		Esta predição é consistente com a Observação~\ref{rem:barreira}:
		assumir $\widehat{w}(0) > 0$ incondicionalmente é equivalente a
		provar Goldbach para $2M^+$. Sob GRH, os erros de Siegel-Walfisz são
		suficientemente controlados para que a cota inferior
		$\widehat{w}(0) > 0$ valha condicionalmente. Esses fatos
		mostram que a reformulação espectral está calibrada corretamente:
		ela captura a mesma dificuldade que a heurística clássica, numa
		linguagem que a expõe mais nitidamente.
	\end{remark}
	
	% ═══════════════════════════════════════════════════════════════════════════════
	\section{Por que os Métodos Clássicos Colapsam}
	\label{sec:crivo}
	
	Esta seção formaliza o diagnóstico mencionado na Introdução,
	descrevendo com precisão as camadas onde a abordagem clássica
	falha. O objetivo não é apenas registrar o colapso, mas mostrar
	por que a formulação espectral evita esses obstáculos específicos
	— ainda que não os resolva.
	
	\subsection{Configuração e notação}
	
	O objeto de estudo nesta seção é o \emph{par alvo} $2M^+$:
	é para ele que $\HRm$ deve ser verificada, pois o par base
	$2M^-$ já é Goldbach por hipótese (dado do motor).
	
	Fixamos $2M^+ = 2M^- + 2$ suficientemente grande e
	$L = \lfloor 2(\log M^+)^2 \rfloor$. Definimos
	\[
	\mathcal{N}(L,\, 2M^+) = \#\{p \leq L : p \in \PP,\;
	2M^+ - p \in \PP\}.
	\]
	A condição $\mathcal{N}(L, 2M^+) \geq 1$ é exatamente $\HRm$
	com a progressão $\{3, 5, \ldots, L\}$,
	$T \sim (\log M^+)^2$.
	
	\subsection{Colapso do método do círculo}
	
	As somas de von Mangoldt
	$f(\alpha) = \sum_{j=1}^T \Lambda(2j+1)\,e(\alpha(2j+1))$
	e $g(\alpha)$ admitem a identidade de Fourier
	\[
	\sum_{j=1}^T \Lambda(2j+1)\,\Lambda(2M^+ - (2j+1))
	= \int_0^1 f(\alpha)\,\overline{g(\alpha)}\,d\alpha.
	\]
	Nos arcos maiores $\mathfrak{M}(a,q)$ com
	$q \leq Q = (\log M^+)^B$, o produto $f\cdot\overline{g}$
	contém o fator $e(2M^+\beta)$ com $|\beta| \leq Q/(qT)$.
	A integral
	\[
	\int_{-Q/T}^{Q/T}|F(\beta,T)|^2\,e(2M^+\beta)\,d\beta
	= O\!\left(\frac{T^2\log T}{M^{+}Q}\right) \to 0
	\]
	quando $M^+ \to \infty$ com $T = O((\log M^+)^2)$ fixo. O
	termo principal é destruído pelas oscilações rápidas de
	$e(2M^+\alpha)$ à escala $M^+ \gg T$.
	
	\subsection{Colapso do crivo de Rosser-Iwaniec}
	
	Crivando $\mathcal{A} = \{2M^+ - p : p \leq L\}$ até nível
	$D = (\log M^+)^{1-\varepsilon}$, o erro acumulado do crivo é
	\[
	\mathcal{R} \leq \sum_{\substack{d \leq D \\ d \mid P(w)}} |r_d|
	\ll D \cdot \frac{L}{(\log L)^A}
	\sim \frac{(\log M^+)^{3-\varepsilon}}{(\log\log M^+)^A},
	\]
	enquanto o termo principal do crivo linear é apenas
	$O((\log M^+)^2/(\log\log M^+)^2)$. Como $3 - \varepsilon > 2$
	para $\varepsilon < 1$, o resto $\mathcal{R}$ domina o sinal e
	a cota inferior colapsa.
	
	\begin{center}
		\renewcommand{\arraystretch}{1.4}
		\begin{tabular}{llll}
			\toprule
			\textbf{Camada} & \textbf{Ferramenta} & \textbf{Resultado}
			& \textbf{Status} \\
			\midrule
			1 & Método do círculo & Integral $\to 0$ para $M^+ \gg T$
			& Colapsa \\
			2 & Bombieri-Vinogradov & Inaplicável: $L \ll M^+$
			& Inaplicável \\
			3 & Crivo linear (Iwaniec) & $S(\mathcal{A},z)$ positivo;
			sozinho não fecha & Parcial \\
			4 & Buchstab + Chen & $\mathcal{R} \gg S(\mathcal{A},z)$
			& Colapsa \\
			\bottomrule
		\end{tabular}
	\end{center}
	
	\subsection{O que a abordagem espectral reposiciona}
	
	A formulação espectral não elimina a dificuldade de provar
	$\widehat{F}(0) > 0$ incondicionalmente. O que ela oferece é:
	\begin{enumerate}[label=(\alph*)]
		\item Um \emph{diagnóstico preciso} da obstrução: a dificuldade
		não está espalhada pelo problema, mas concentrada na componente
		de frequência zero de um único objeto espectral.
		\item Uma \emph{linguagem alinhada} com as técnicas que
		resolveram problemas análogos: Green-Tao usou exatamente a
		transição ``contagem $\to$ fluxo'' para progressões aritméticas
		em primos; Helfgott \cite{helfgott} usou análise de Fourier
		refinada para Goldbach ternário.
		\item Uma \emph{separação clara} entre o que está provado (a
		estrutura dinâmica do motor, incluindo os quatro eixos de $W_i$
		e o mecanismo dos pivôs do Scanner) e o que permanece aberto
		(a positividade de $\widehat{F}(0)$).
	\end{enumerate}
	
	% ═══════════════════════════════════════════════════════════════════════════════
	\section{Síntese, Alcance e Trabalhos Futuros}
	\label{sec:conclusao}
	
	\subsection{O que este artigo prova}
	
	\begin{center}
		\renewcommand{\arraystretch}{1.4}
		\begin{tabular}{lll}
			\toprule
			\textbf{Resultado} & \textbf{Status} & \textbf{Referência} \\
			\midrule
			Fita-Dobra: colunas $= 2M^-$, diagonais $= 2M^+$
			& Incondicional & Def.~\ref{def:fita},
			Props.~\ref{prop:coluna}--\ref{prop:diagonal} \\
			Grade zigzag e acoplamento $3C = 2N$
			& Incondicional & Def.~\ref{def:grade}--\ref{def:acoplamento} \\
			$W_i$: 4 eixos, soma constante $= 2M^+$
			& Incondicional & Def.~\ref{def:wi}, Prop.~\ref{prop:eixos} \\
			Scanner com pivôs, cobertura de $N-1$ diagonais
			& Incondicional & Def.~\ref{def:scanner} \\
			Sistema $(X_N,\mu_N,T)$ bem definido
			& Incondicional & Def.~\ref{def:estados}--\ref{def:koopman} \\
			$T$ ergódico, $h(T) = 0$
			& Incondicional & Lema~\ref{lem:ergodico} \\
			Espectro discreto de $U_T$ (raízes da unidade)
			& Incondicional & Teo.~\ref{thm:espectro} \\
			$\HRm \Leftrightarrow \widehat{F}(0) > 0$
			& Incondicional & Teo.~\ref{thm:reformulacao} \\
			Descorrelação $\widehat{w}(j) \to 0$ para $j \neq 0$
			& Incondicional (GTZ) & Prop.~\ref{prop:sarnak} \\
			$\widehat{w}(0) > 0$
			& Cond.\ (H-L ou GRH) & Obs.~\ref{rem:HL} \\
			\midrule
			\textbf{$\HRm$ provada} & \textbf{Não provada}
			& Obs.~\ref{rem:barreira} \\
			\bottomrule
		\end{tabular}
	\end{center}
	
	A última linha da tabela é o ponto central: o artigo reformula
	$\HRm$, não a demonstra. A equivalência
	$\HRm \Leftrightarrow \widehat{F}(0) > 0$ é uma tautologia
	precisa cuja utilidade está em tornar a barreira analítica
	explícita e acessível às ferramentas modernas.
	
	\subsection{O que a reformulação oferece à comunidade}
	
	\begin{enumerate}
		\item \textbf{Localização da obstrução.} A Conjectura de
		Goldbach para $2M^+$ equivale à positividade de uma única
		componente de Fourier $\widehat{F}(0)$. Isso não simplifica
		o problema, mas o \emph{concentra}: em vez de uma afirmação
		difusa sobre todos os pares, temos uma afirmação espectral
		localizada.
		
		\item \textbf{Alinhamento com a escola moderna.} O Motor de
		Herança, formulado como sistema de Koopman de entropia zero,
		encaixa-se naturalmente na estrutura onde o teorema de
		nilsequências de Green-Tao-Ziegler e a Conjectura de Sarnak são
		aplicáveis. A descorrelação das componentes $j \neq 0$ já é
		incondicional (Proposição~\ref{prop:sarnak}); a única barreira
		remanescente é a componente de frequência zero.
		
		\item \textbf{Separação limpa entre camadas.} A camada algébrica
		(incondicional) — que inclui agora a estrutura completa de
		$W_i$ com quatro eixos e o mecanismo dos pivôs do Scanner —
		e a camada aritmética (condicional) estão separadas com
		precisão. Qualquer progresso futuro em $\widehat{F}(0)$ se
		beneficia de toda a estrutura dinâmica já estabelecida.
	\end{enumerate}
	
	\subsection{Direções futuras}
	
	\begin{enumerate}
		\item \textbf{Nilsequências e normas de Gowers de ordem
			superior.} A decomposição de $F$ em contribuições de
		nilsequências (no espírito de Green-Tao-Ziegler \cite{greentao})
		pode fornecer controle sobre $\widehat{F}(0)$, se a progressão
		$\{E_1(t)\}_{t=1}^K$ exibir estrutura de nilsequência.
		
		\item \textbf{Prova condicional sob GRH.} Sob a Hipótese de
		Riemann Generalizada, os erros de Siegel-Walfisz satisfazem
		$|r_d| \ll d^{1/2}(\log M^+)^2$, e
		$\mathcal{R} \ll D^{3/2}(\log M^+)^2$. Para
		$D = (\log M^+)^{1-\varepsilon}$ com $\varepsilon$ próximo de
		$1$, isso tornaria $\mathcal{R} \ll S(\mathcal{A},w)$ e fecharia
		a cota inferior $\mathcal{N} \geq 1$ condicionalmente.
		
		\item \textbf{Recorrência de Furstenberg.} A ergodicidade de $T$
		e a densidade positiva de $A_N$ (sob hipótese de
		Hardy-Littlewood) sugerem uma prova de recorrência múltipla: a
		órbita de qualquer ponto em $X_N$ deveria visitar $A_N$ no
		horizonte $K = N-1$. Formalizar isso via teoremas de recorrência
		de Furstenberg-Katznelson é uma direção concreta.
		
		\item \textbf{Análise do Tempo de Primeira Passagem.} O
		experimento computacional de transição Caos $\to$ Ordem
		(\cite{motor}, Seção~8) mede quantos passos do Scanner até que
		$\HRm$ seja ativada. A análise assintótica desse tempo em função
		de $N$ — e sua conexão com a Lei de Saturação Geométrica —
		é uma direção concreta para verificação computacional de
		$\widehat{F}(0) > 0$ em escala.
	\end{enumerate}
	
	% ── Referências ───────────────────────────────────────────────────────────────
	\begin{thebibliography}{99}
		
		\bibitem{goldbach1}
		T.~Bandeira,
		\textit{Uma Abordagem Unificada para a Conjectura de Goldbach:
			Estrutura Combinatória, Saturação Geométrica e o Elo entre
			Trios e Pares Primos},
		Preprint, maio de 2026.
		
		\bibitem{grade}
		T.~Bandeira,
		\textit{Uma propriedade de soma constante na grade $3 \times N$
			e as conjecturas de Goldbach},
		Preprint, maio de 2026.
		
		\bibitem{ancora}
		T.~Bandeira,
		\textit{O Invariante Âncora da Grade $G_{3,N}$ e um Estimador
			Geométrico Oracle-Free para a Conjectura de Goldbach},
		Preprint, maio de 2026.
		
		\bibitem{crivo_b}
		T.~Bandeira,
		\textit{Problema B do Crivo Geométrico: Demonstração da
			Desigualdade Estrutural},
		Preprint, maio de 2026.
		
		\bibitem{motor}
		T.~Bandeira,
		\textit{Motor de Herança Estrutural: Formalização Definitiva
			da Fita-Dobra, da Grade Zigzag, da Janela $W_i$ com Quatro
			Eixos e do Mecanismo de Pivôs},
		Preprint, maio de 2026.
		
		\bibitem{greentao}
		B.~Green e T.~Tao,
		\textit{The primes contain arbitrarily long arithmetic
			progressions},
		Annals of Mathematics, vol.~167, pp.~481--547, 2008.
		
		\bibitem{greentaoziegler}
		B.~Green, T.~Tao e T.~Ziegler,
		\textit{An inverse theorem for the Gowers $U^{s+1}[N]$-norm},
		Annals of Mathematics, vol.~176, pp.~1231--1372, 2012.
		
		\bibitem{gowers}
		W.~T.~Gowers,
		\textit{A new proof of Szemerédi's theorem},
		Geometric and Functional Analysis, vol.~11, pp.~465--588, 2001.
		
		\bibitem{sarnak}
		P.~Sarnak,
		\textit{Three lectures on the Möbius function, randomness and
			dynamics},
		Publications mathématiques de l'IHÉS, 2011.
		
		\bibitem{helfgott}
		H.~A.~Helfgott,
		\textit{Major arcs for Goldbach's theorem},
		arXiv:1305.2897, 2013.
		
		\bibitem{hardy-littlewood}
		G.~H.~Hardy e J.~E.~Littlewood,
		\textit{Some problems of `Partitio Numerorum'; III: On the
			expression of a number as a sum of primes},
		Acta Mathematica, vol.~44, pp.~1--70, 1923.
		
		\bibitem{iwaniec}
		H.~Iwaniec e E.~Kowalski,
		\textit{Analytic Number Theory},
		American Mathematical Society Colloquium Publications,
		vol.~53, 2004.
		
		\bibitem{chen}
		J.~R.~Chen,
		\textit{On the representation of a large even integer as the
			sum of a prime and the product of at most two primes},
		Scientia Sinica, vol.~16, pp.~157--176, 1973.
		
	\end{thebibliography}
	
\end{document}